\chapter{Impfung und Schwellenwerte}

\section{Die Frage}

\begin{merksatz}
\boxtitel{Neustadt, ein Jahr später}
Gegen den Erreger steht nun ein Impfstoff zur Verfügung. Die Gesundheitsbehörde fragt:

\begin{center}
\begin{minipage}{0.9\textwidth}
\textbf{Wie viele Menschen müssen geimpft werden, damit sich der Erreger nicht mehr ausbreiten kann?}
\end{minipage}
\end{center}

Die Antwort lautet nicht \enquote{alle} --- und genau darin liegt der Wert des Modells.
\end{merksatz}

\begin{merksatz}
\boxtitel{Satz 8.1 --- die Herdenschwelle}
Wird ein Anteil $p$ der Bevölkerung vor Beginn immunisiert, so breitet sich der Erreger genau dann nicht aus, wenn
\[ p \;>\; p_c \;:=\; 1 - \frac{1}{\Rnull} . \]

\schritt{Beweis}
Geimpfte gehören von Anfang an zur Kammer $R$. Der Anteil der Empfänglichen beträgt also $S/N = 1-p$.

Nach Satz 7.2 breitet sich der Erreger aus, wenn
\[ \Rnull\cdot\frac{S}{N} > 1 \quad\Longleftrightarrow\quad \Rnull\,(1-p) > 1 . \]

Verneint man das:
\[ \Rnull\,(1-p) < 1 \quad\Longleftrightarrow\quad 1-p < \frac{1}{\Rnull} \quad\Longleftrightarrow\quad p > 1 - \frac{1}{\Rnull} . \qquad \blacksquare \]
\end{merksatz}

\begin{merksatz}
\boxtitel{Warum nicht alle geimpft werden müssen}
\begin{center}
\begin{minipage}{0.9\textwidth}
Eine ansteckende Person trifft im Mittel $\Rnull$ empfängliche Menschen --- \kursiv{wenn} alle empfänglich sind. Ist nur der Anteil $1-p$ empfänglich, so sind es nur $\Rnull(1-p)$.
\end{minipage}
\end{center}

Man muss also nicht jede Ansteckung verhindern, sondern nur so viele, dass jeder Ansteckende im Mittel weniger als einen neuen erzeugt.

\begin{center}
\begin{tikzpicture}[scale=1.0]
\begin{scope}
\foreach \i in {0,...,24} {
  \pgfmathsetmacro{\xx}{mod(\i,5)*0.42}
  \pgfmathsetmacro{\yy}{-floor(\i/5)*0.42}
  \fill[grau!40] (\xx,\yy) circle (2.4pt);
}
\fill[akzent] (0.84,-0.84) circle (3.2pt);
\node[font=\scriptsize\itshape, color=grau] at (0.84,-2.4) {ohne Impfung};
\end{scope}
\begin{scope}[xshift=4.6cm]
\foreach \i in {0,...,24} {
  \pgfmathsetmacro{\xx}{mod(\i,5)*0.42}
  \pgfmathsetmacro{\yy}{-floor(\i/5)*0.42}
  \fill[grau!40] (\xx,\yy) circle (2.4pt);
}
\foreach \i in {0,2,3,5,7,8,10,11,13,15,16,18,20,21,23} {
  \pgfmathsetmacro{\xx}{mod(\i,5)*0.42}
  \pgfmathsetmacro{\yy}{-floor(\i/5)*0.42}
  \fill[gruen] (\xx,\yy) circle (2.4pt);
}
\fill[akzent] (0.84,-0.84) circle (3.2pt);
\node[font=\scriptsize\itshape, color=grau] at (0.84,-2.4) {mit $60\,\%$ geimpft};
\end{scope}
\node[font=\scriptsize, color=grau, anchor=west] at (7.2,-0.4) {grün: geschützt};
\node[font=\scriptsize, color=akzent, anchor=west] at (7.2,-0.9) {rot: ansteckend};
\end{tikzpicture}
\end{center}
\end{merksatz}

\begin{merksatz}
\boxtitel{Die Schwellen für bekannte Erreger}
\begin{center}
{\small
\renewcommand{\arraystretch}{1.48}
\begin{tabular}{@{}P{2.6cm}P{2.2cm}P{2.6cm}@{}}
\toprule
\textbf{Erreger} & $\Rnull$ & \textbf{Schwelle $p_c$}\\
\midrule
Masern & $12$--$18$ & $92$--$94\,\%$\\
Keuchhusten & $12$--$17$ & $92$--$94\,\%$\\
Windpocken & $8$--$12$ & $88$--$92\,\%$\\
Röteln & $6$--$7$ & $83$--$86\,\%$\\
Pocken & $5$--$7$ & $80$--$86\,\%$\\
Mumps & $4$--$7$ & $75$--$86\,\%$\\
Neustädter Erreger & $2{,}5$ & $\mathbf{60\,\%}$\\
Grippe, saisonal & $1{,}2$--$1{,}4$ & $17$--$29\,\%$\\
\bottomrule
\end{tabular}}
\end{center}

\begin{center}
\begin{minipage}{0.88\textwidth}
\kursiv{Die Spannen sind echt: $\Rnull$ hängt von Bevölkerungsdichte, Altersaufbau und Verhalten ab und ist für denselben Erreger in verschiedenen Gesellschaften verschieden. Die Werte für Masern gehören zu den höchsten, die man kennt --- deshalb sind Masern der schwierigste Fall und zugleich der empfindlichste Anzeiger für Lücken im Impfschutz.}
\end{minipage}
\end{center}

\schritt{Die Antwort für Neustadt}
Bei $\Rnull = 2{,}5$ genügen $60\,\%$. Vierzig Prozent der Bevölkerung müssen nicht geimpft werden --- und sind trotzdem geschützt.
\end{merksatz}

\section{Unvollständiger Schutz}

\begin{warnung}
\boxtitelr{Kein Impfstoff wirkt bei allen}
Sei $\varepsilon$ der Anteil der Geimpften, bei denen der Impfschutz greift --- die \Begriff{Wirksamkeit}.

\schritt{Die neue Bedingung}
Der Anteil tatsächlich Geschützter beträgt $p\,\varepsilon$. Also
\[ p\,\varepsilon > 1 - \frac{1}{\Rnull} \quad\Longleftrightarrow\quad \boxed{\;p > \frac{1 - 1/\Rnull}{\varepsilon}\;} \]

\schritt{Die unbequeme Folgerung}
\begin{center}
\begin{minipage}{0.9\textwidth}
\textbf{Ist $\varepsilon < 1 - 1/\Rnull$, so lässt sich die Schwelle auch dann nicht erreichen, wenn jeder einzelne Mensch geimpft wird.}
\end{minipage}
\end{center}

Denn dann ist $p > 1$ verlangt, was unmöglich ist.

\schritt{Am Beispiel}
\begin{center}
{\small
\renewcommand{\arraystretch}{1.48}
\begin{tabular}{@{}P{1.8cm}P{1.6cm}P{2.2cm}P{2.4cm}@{}}
\toprule
$\Rnull$ & $p_c$ & $\varepsilon = 0{,}95$ & $\varepsilon = 0{,}60$\\
\midrule
$2{,}5$ & $60{,}0\,\%$ & $63{,}2\,\%$ & $100{,}0\,\%$\\
$6$ & $83{,}3\,\%$ & $87{,}7\,\%$ & \textcolor{akzent}{unmöglich}\\
$15$ & $93{,}3\,\%$ & $98{,}2\,\%$ & \textcolor{akzent}{unmöglich}\\
\bottomrule
\end{tabular}}
\end{center}

\kursiv{Rechnung für $\Rnull=15$, $\varepsilon=0{,}95$:}
\[ p > \frac{0{,}9333}{0{,}95} = 0{,}9824 \]

\begin{center}
\begin{minipage}{0.88\textwidth}
\kursiv{Bei Masern muss also praktisch jeder geimpft werden, obwohl der Impfstoff zu den wirksamsten gehört. Die zweite Impfung, die in vielen Ländern empfohlen wird, dient genau diesem Zweck: Sie hebt die Wirksamkeit von etwa $93\,\%$ auf über $97\,\%$ und macht die Schwelle dadurch überhaupt erst erreichbar.}
\end{minipage}
\end{center}
\end{warnung}

\section{Wenn die Schwelle nicht erreicht wird}

Der wichtigere Fall in der Praxis: Was nützt eine Impfquote \kursiv{unterhalb} der Schwelle?

\begin{merksatz}
\boxtitel{Die Endgröße bei teilweiser Impfung}
Wir wenden Satz 7.4 an, aber mit dem Ausgangsanteil $s_0 = 1-p$ statt $1$. Dieselbe Herleitung ergibt
\[ s_\infty = s_0\,e^{-\Rnull\,(s_0 - s_\infty)} . \]

\schritt{Für Neustadt, $\Rnull = 2{,}5$}
\begin{center}
{\small
\renewcommand{\arraystretch}{1.48}
\begin{tabular}{@{}P{1.4cm}P{1.6cm}P{2.2cm}P{2.6cm}@{}}
\toprule
$p$ & $s_0$ & \textbf{Erkrankte} & \textbf{davon bei den Ungeimpften}\\
\midrule
$0$ & $1{,}00$ & $89{,}3\,\%$ & $89{,}3\,\%$\\
$0{,}10$ & $0{,}90$ & $76{,}8\,\%$ & $85{,}3\,\%$\\
$0{,}20$ & $0{,}80$ & $63{,}7\,\%$ & $79{,}7\,\%$\\
$0{,}30$ & $0{,}70$ & $49{,}9\,\%$ & $71{,}3\,\%$\\
$0{,}40$ & $0{,}60$ & $35{,}0\,\%$ & $58{,}3\,\%$\\
$0{,}50$ & $0{,}50$ & $18{,}6\,\%$ & $37{,}1\,\%$\\
$0{,}55$ & $0{,}45$ & $9{,}6\,\%$ & $21{,}4\,\%$\\
$\mathbf{0{,}60}$ & $0{,}40$ & $\mathbf{0{,}2\,\%}$ & $0{,}4\,\%$\\
\bottomrule
\end{tabular}}
\end{center}

\schritt{Die Deutung}
\begin{center}
\begin{minipage}{0.9\textwidth}
\textbf{Jede Impfung nützt --- und zwar überproportional.} Zehn Prozent Impfquote senken die Zahl der Erkrankten um $12{,}5$ Prozentpunkte, nicht um zehn.
\end{minipage}
\end{center}

\schritt{Der indirekte Schutz}
Die rechte Spalte zeigt den Anteil der \kursiv{Ungeimpften}, die erkranken. Er fällt von $89{,}3$ auf $21{,}4$ Prozent --- ohne dass diese Menschen etwas getan hätten.

\begin{center}
\begin{minipage}{0.88\textwidth}
\kursiv{Das ist der \Begriff{indirekte Schutz}, und er ist der eigentliche Gegenstand dieses Kapitels. Eine Impfung nützt nicht nur der geimpften Person; sie nimmt den Erreger aus dem Umlauf und schützt dadurch auch diejenigen, die nicht geimpft werden können --- Säuglinge, Menschen mit geschwächtem Abwehrsystem, Personen mit Gegenanzeigen.}
\end{minipage}
\end{center}

\schritt{Die Nichtlinearität am Rand}
Von $p = 0{,}55$ auf $p = 0{,}60$ --- fünf Prozentpunkte --- fällt die Erkranktenzahl von $9{,}6$ auf $0{,}2$ Prozent.

\begin{center}
\begin{minipage}{0.9\textwidth}
\textbf{Die letzten Prozentpunkte vor der Schwelle bewirken am meisten.} Genau dort ist die Impfbereitschaft erfahrungsgemäß am schwersten zu erhöhen.
\end{minipage}
\end{center}
\end{merksatz}

\section{Die Altersverschiebung}

Nun ein Ergebnis, das der Anschauung zuwiderläuft und praktisch bedeutsam ist.

\begin{merksatz}
\boxtitel{Das Modell mit Nachwuchs}
Für Erreger, die über Jahrzehnte umlaufen, muss man Geburten berücksichtigen:
\[ \frac{\dd S}{\dd t} = \mu N - \beta\frac{SI}{N} - \mu S, \qquad
\frac{\dd I}{\dd t} = \beta\frac{SI}{N} - \gamma I - \mu I \]
mit der Geburts- und Sterberate $\mu = 1/L$, wobei $L$ die mittlere Lebensdauer ist.

\schritt{Das endemische Gleichgewicht}
Aus $\dd I/\dd t = 0$ mit $I \neq 0$:
\[ \beta\frac{S}{N} = \gamma + \mu \quad\Longrightarrow\quad \frac{S^{*}}{N} = \frac{\gamma+\mu}{\beta} = \frac{1}{\Rnull} \]
mit $\Rnull = \beta/(\gamma+\mu)$.

\begin{center}
\begin{minipage}{0.9\textwidth}
\textbf{Im Dauerzustand ist genau der Anteil $1/\Rnull$ empfänglich.} Die Krankheit reguliert sich selbst auf die Schwelle ein.
\end{minipage}
\end{center}

\begin{center}
\begin{minipage}{0.88\textwidth}
\kursiv{Das ist bemerkenswert: Der Anteil hängt nur von $\Rnull$ ab, nicht von der Geburtenrate. Und es liefert ein Verfahren zur Schätzung von $\Rnull$ aus Antikörperuntersuchungen --- misst man, welcher Anteil einer Bevölkerung noch nie Kontakt hatte, so kennt man $1/\Rnull$.}
\end{minipage}
\end{center}
\end{merksatz}

\begin{merksatz}
\boxtitel{Satz 8.2 --- das mittlere Erkrankungsalter}
Im endemischen Gleichgewicht beträgt das mittlere Alter bei Ansteckung näherungsweise
\[ A \;\approx\; \frac{L}{\Rnull - 1} . \]

\schritt{Begründung}
Eine empfängliche Person wird mit der Rate $\lambda = \beta I^{*}/N$ angesteckt. Die mittlere Wartezeit bis dahin ist $1/\lambda$ --- wie bei jedem exponentiellen Vorgang (Kapitel~1).

Aus dem Gleichgewicht der ersten Gleichung ergibt sich nach kurzer Rechnung $\lambda = \mu(\Rnull-1)$, also
\[ A = \frac{1}{\mu(\Rnull-1)} = \frac{L}{\Rnull-1} . \]

\schritt{Am Beispiel Masern}
$\Rnull = 15$, $L = 80$ Jahre:
\[ A = \frac{80}{14} = 5{,}7 \ \text{Jahre} \]

\begin{center}
\begin{minipage}{0.88\textwidth}
\kursiv{Das stimmt mit der Beobachtung überein: Vor Einführung der Impfung war Masern eine Krankheit des Vorschulalters. Je ansteckender ein Erreger, desto früher trifft er --- weil er praktisch jeden trifft, und zwar so bald wie möglich.}
\end{minipage}
\end{center}
\end{merksatz}

\begin{warnung}
\boxtitelr{Was die Impfung mit dem Erkrankungsalter macht}
Bei einer Impfquote $p$ unterhalb der Schwelle sinkt die effektive Reproduktionszahl auf $\Rnull(1-p)$. Also:
\[ A_p \approx \frac{L}{\Rnull(1-p) - 1} \]

\schritt{Für Röteln --- $\Rnull = 6$, $L = 75$ Jahre}
\begin{center}
{\small
\renewcommand{\arraystretch}{1.48}
\begin{tabular}{@{}P{1.6cm}P{2.0cm}P{2.6cm}@{}}
\toprule
$p$ & $\Rnull(1-p)$ & \textbf{mittleres Alter}\\
\midrule
$0$ & $6{,}00$ & $15{,}0$ Jahre\\
$0{,}30$ & $4{,}20$ & $23{,}4$\\
$0{,}50$ & $3{,}00$ & \textcolor{akzent}{$\mathbf{37{,}5}$}\\
$0{,}70$ & $1{,}80$ & $93{,}7$\\
$\mathbf{0{,}833}$ & $1{,}00$ & \textcolor{gruen}{Elimination}\\
\bottomrule
\end{tabular}}
\end{center}

\begin{center}
\begin{minipage}{0.9\textwidth}
\textbf{Eine Impfquote von $50\,\%$ verschiebt das mittlere Erkrankungsalter von $15$ auf $37$ Jahre.}
\end{minipage}
\end{center}

\schritt{Warum das gefährlich ist}
Röteln sind im Kindesalter harmlos. In der Schwangerschaft führen sie zu schweren Fehlbildungen beim Kind.

\begin{center}
\begin{minipage}{0.88\textwidth}
\kursiv{Eine unzureichende Impfquote kann daher die Zahl der schweren Fälle \kursiv{erhöhen}, obwohl die Gesamtzahl der Erkrankungen sinkt. Genau das geschah in Griechenland in den 1990er Jahren: Nach jahrelang mäßiger Impfquote kam es 1993 zu einem Rötelnausbruch mit ungewöhnlich vielen Erkrankungen im gebärfähigen Alter und in der Folge zur höchsten je dort verzeichneten Zahl angeborener Rötelnschäden.}
\end{minipage}
\end{center}

\schritt{Die Folgerung}
\begin{center}
\begin{minipage}{0.9\textwidth}
\textbf{Bei Erregern, deren Schwere mit dem Alter zunimmt, ist eine halbherzige Impfkampagne schlechter als keine.}
\end{minipage}
\end{center}

Das betrifft Röteln, Mumps, Windpocken, Hepatitis~A und Polio. Es betrifft \kursiv{nicht} Masern, deren Verlauf im Säuglingsalter am schwersten ist.

\begin{center}
\begin{minipage}{0.88\textwidth}
\kursiv{Die praktische Folgerung, die daraus gezogen wurde: Wer gegen Röteln impft, muss die Schwelle erreichen --- oder gezielt diejenigen impfen, für die die Krankheit gefährlich ist. Beides ist heute Bestandteil der Empfehlungen.}
\end{minipage}
\end{center}
\end{warnung}

\section{Der Bruchfall}

\begin{warnung}
\boxtitelr{Warum die Schwelle in der Wirklichkeit anders liegt}
\begin{center}
{\small
\renewcommand{\arraystretch}{1.5}
\begin{tabular}{@{}P{2.4cm}P{2.8cm}P{2.8cm}@{}}
\toprule
\textbf{Abweichung} & \textbf{Ursache} & \textbf{Folge}\\
\midrule
Heterogenität & Kontaktzahlen streuen & \textcolor{gruen}{Schwelle sinkt}\\
Bündelung & Ungeimpfte finden einander & \textcolor{akzent}{Schwelle steigt}\\
Nachlassender Schutz & Immunität schwindet & \textcolor{akzent}{Auffrischung nötig}\\
Einschleppung & keine geschlossene Bevölkerung & Ausbrüche bleiben möglich\\
\bottomrule
\end{tabular}}
\end{center}

\schritt{Die zweite Zeile ist die folgenreichste}
Das Modell setzt voraus, dass Geimpfte und Ungeimpfte gleichmäßig gemischt sind. Tatsächlich sammeln sich Impfgegner in bestimmten Schulen, Gemeinden und Regionen.

\begin{center}
\begin{minipage}{0.9\textwidth}
\textbf{Eine landesweite Impfquote von $95\,\%$ schützt nicht vor Ausbrüchen, wenn es einzelne Gemeinschaften mit $50\,\%$ gibt.} Dort ist die örtliche Reproduktionszahl über eins.
\end{minipage}
\end{center}

\begin{center}
\begin{minipage}{0.88\textwidth}
\kursiv{Alle größeren Masernausbrüche der letzten Jahrzehnte in Ländern mit hoher Durchimpfung folgten diesem Muster: Sie begannen in räumlich oder sozial abgegrenzten Gruppen mit niedriger Quote. Die Durchschnittszahl verdeckt das --- deshalb ist die \kursiv{Verteilung} der Impfquote wichtiger als ihr Mittelwert.}
\end{minipage}
\end{center}

\schritt{Zur ersten Zeile}
Streuen die Kontaktzahlen stark, so genügt es, die Kontaktfreudigsten zu impfen. Die Schwelle sinkt dann unter $1-1/\Rnull$.

\begin{center}
\begin{minipage}{0.88\textwidth}
\kursiv{Das ist die Grundlage gezielter Impfstrategien --- etwa der Riegelungsimpfung im Umfeld eines Falls, mit der die Pocken ausgerottet wurden. Sie wirkt, weil sie genau die Personen erreicht, über die die Übertragung läuft.}
\end{minipage}
\end{center}
\end{warnung}

\section{Übungsaufgaben}

\begin{aufgabe}
\boxtitel{Aufgabe 8.1 --- Schwelle mit unvollständiger Wirksamkeit}
Für einen Erreger mit $\Rnull = 8$ steht ein Impfstoff mit $\varepsilon = 0{,}90$ zur Verfügung. Welche Quote ist nötig? Was ändert sich bei $\varepsilon = 0{,}85$?

\loesung
\schritt{Schritt 1 --- Die Grundschwelle}
\[ p_c = 1 - \frac{1}{8} = 0{,}875 = 87{,}5\,\% \]

\schritt{Schritt 2 --- Mit $\varepsilon = 0{,}90$}
\[ p > \frac{0{,}875}{0{,}90} = 0{,}9722 = 97{,}2\,\% \]

\schritt{Schritt 3 --- Mit $\varepsilon = 0{,}85$}
\[ p > \frac{0{,}875}{0{,}85} = 1{,}0294 \]

Über eins --- \textbf{unmöglich.} Selbst bei vollständiger Durchimpfung bliebe die effektive Reproduktionszahl über eins.

\schritt{Schritt 4 --- Die Probe}
Bei $p=1$ und $\varepsilon = 0{,}85$ sind $85\,\%$ geschützt, $15\,\%$ empfänglich:
\[ \mathcal{R}_{\text{eff}} = 8 \cdot 0{,}15 = 1{,}20 > 1 \quad \textcolor{akzent}{\times} \]
Bei $\varepsilon = 0{,}90$: $\mathcal{R}_{\text{eff}} = 8 \cdot 0{,}10 = 0{,}80 < 1$ \checkmark

\schritt{Schritt 5 --- Die kritische Wirksamkeit}
Allgemein muss gelten $\varepsilon \ge p_c$, also hier $\varepsilon \ge 0{,}875$.

\begin{center}
\begin{minipage}{0.9\textwidth}
\textbf{Die Wirksamkeit muss mindestens so groß sein wie die Herdenschwelle.} Andernfalls ist Elimination durch Impfung allein ausgeschlossen.
\end{minipage}
\end{center}

\schritt{Schritt 6 --- Die Empfindlichkeit}
\begin{center}
\renewcommand{\arraystretch}{1.45}
\begin{tabular}{@{}P{2.0cm}P{2.4cm}P{2.4cm}@{}}
\toprule
$\varepsilon$ & \textbf{nötige Quote} & \\
\midrule
$1{,}00$ & $87{,}5\,\%$ & \\
$0{,}95$ & $92{,}1\,\%$ & \\
$0{,}90$ & $97{,}2\,\%$ & \\
$0{,}88$ & $99{,}4\,\%$ & \textcolor{akzent}{kaum erreichbar}\\
$0{,}875$ & $100{,}0\,\%$ & \textcolor{akzent}{Grenze}\\
\bottomrule
\end{tabular}
\end{center}

\begin{center}
\begin{minipage}{0.86\textwidth}
\kursiv{Fünf Prozentpunkte Wirksamkeit --- von $95$ auf $90$ --- erfordern fünf Prozentpunkte mehr Impfquote. Die letzten Zeilen zeigen, wie steil es danach wird: Nahe der kritischen Wirksamkeit wächst die nötige Quote über jede Grenze.}
\end{minipage}
\end{center}
\end{aufgabe}

\begin{aufgabe}
\boxtitel{Aufgabe 8.2 --- Die Altersverschiebung beziffern}
Für Mumps gilt $\Rnull = 5$ und $L = 78$ Jahre. Bestimmen Sie das mittlere Erkrankungsalter ohne Impfung und bei Quoten von $40$, $60$ und $75$ Prozent. Ab welcher Quote läge es über $20$ Jahren?

\loesung
\schritt{Schritt 1 --- Ohne Impfung}
\[ A = \frac{78}{5-1} = \frac{78}{4} = 19{,}5 \ \text{Jahre} \]

\begin{center}
\begin{minipage}{0.86\textwidth}
\kursiv{Schon ohne Impfung liegt das mittlere Alter bei knapp zwanzig Jahren. Mumps ist damit von vornherein ein Grenzfall --- die gefürchtete Hodenentzündung tritt vor allem nach der Pubertät auf.}
\end{minipage}
\end{center}

\schritt{Schritt 2 --- Die Schwelle}
\[ p_c = 1 - \frac{1}{5} = 0{,}80 = 80\,\% \]

\schritt{Schritt 3 --- Die Werte}
\begin{center}
\renewcommand{\arraystretch}{1.48}
\begin{tabular}{@{}P{1.6cm}P{2.2cm}P{2.6cm}@{}}
\toprule
$p$ & $\Rnull(1-p)$ & $A_p$\\
\midrule
$0$ & $5{,}00$ & $19{,}5$ J\\
$0{,}40$ & $3{,}00$ & $39{,}0$ J\\
$0{,}60$ & $2{,}00$ & $78{,}0$ J\\
$0{,}75$ & $1{,}25$ & $312{,}0$ J\\
$0{,}80$ & $1{,}00$ & \textcolor{gruen}{Elimination}\\
\bottomrule
\end{tabular}
\end{center}

\kursiv{Rechnung für $p=0{,}40$:} $A = 78/(3-1) = 39{,}0$.

\schritt{Schritt 4 --- Die Deutung der großen Werte}
Ein \enquote{mittleres Erkrankungsalter} von $312$ Jahren ist natürlich unsinnig. Die Zahl bedeutet: Die Ansteckungsrate ist so gering, dass die allermeisten Menschen \kursiv{nie} erkranken --- sie sterben vorher.

\begin{center}
\begin{minipage}{0.86\textwidth}
\kursiv{Die Formel gilt nur, solange $A$ klein gegen $L$ ist. Bei $A \ge L$ ist sie als Größenvergleich zu lesen, nicht als Vorhersage. Das ist eine Bruchstelle des Modells, die man kennen muss --- die Rechnung liefert weiter Zahlen, aber sie bedeuten etwas anderes.}
\end{minipage}
\end{center}

\schritt{Schritt 5 --- Ab welcher Quote über zwanzig Jahren}
\[ \frac{78}{5(1-p)-1} > 20 \]
\[ 5(1-p) - 1 < \frac{78}{20} = 3{,}9 \]
\[ 5(1-p) < 4{,}9 \quad\Longrightarrow\quad 1-p < 0{,}98 \quad\Longrightarrow\quad p > 0{,}02 \]

\schritt{Antwort}
\textbf{Ab zwei Prozent Impfquote.}

\begin{center}
\begin{minipage}{0.9\textwidth}
\textbf{Bei Mumps genügt schon eine minimale Impfquote, um das mittlere Erkrankungsalter über zwanzig Jahre zu schieben} --- weil es ohne Impfung bereits bei $19{,}5$ lag.
\end{minipage}
\end{center}

\begin{center}
\begin{minipage}{0.88\textwidth}
\kursiv{Das ist die schärfste Form des Arguments aus Abschnitt 8.4: Bei einem Erreger, dessen Erkrankungsalter ohnehin nahe der kritischen Grenze liegt, kippt jede unvollständige Impfung die Lage ins Ungünstige. Die einzige vertretbare Strategie ist dann, die Schwelle zu erreichen.}
\end{minipage}
\end{center}
\end{aufgabe}

\begin{aufgabe}
\boxtitel{Aufgabe 8.3 --- Wem nützt die Impfung?}
In Neustadt sind $50\,\%$ geimpft, $\Rnull = 2{,}5$. Von den Ungeimpften erkranken $37{,}1\,\%$. Vergleichen Sie den direkten und den indirekten Nutzen.

\loesung
\schritt{Schritt 1 --- Ohne jede Impfung}
Alle $100\,000$ sind empfänglich, $89\,300$ erkranken.

\schritt{Schritt 2 --- Mit $50\,\%$ Impfquote}
$50\,000$ Geimpfte erkranken nicht. Von den $50\,000$ Ungeimpften erkranken $37{,}1\,\%$:
\[ 0{,}371 \cdot 50\,000 = 18\,550 \]

\schritt{Schritt 3 --- Die Bilanz}
\begin{center}
\renewcommand{\arraystretch}{1.48}
\begin{tabular}{@{}P{3.0cm}P{2.2cm}P{2.4cm}@{}}
\toprule
& \textbf{ohne} & \textbf{mit $p=0{,}5$}\\
\midrule
Erkrankte gesamt & $89\,300$ & $18\,550$\\
Vermieden & --- & $\mathbf{70\,750}$\\
\bottomrule
\end{tabular}
\end{center}

\schritt{Schritt 4 --- Wer profitiert wovon}
\kursiv{Direkter Nutzen:} Die $50\,000$ Geimpften hätten ohne Impfung zu $89{,}3\,\%$ erkrankt, also
\[ 0{,}893 \cdot 50\,000 = 44\,650 \ \text{Fälle} . \]

\kursiv{Indirekter Nutzen:} Die $50\,000$ Ungeimpften hätten $44\,650$ Fälle gehabt, haben aber nur $18\,550$:
\[ 44\,650 - 18\,550 = 26\,100 \ \text{vermiedene Fälle} . \]

\begin{center}
\renewcommand{\arraystretch}{1.48}
\begin{tabular}{@{}P{3.0cm}P{2.2cm}P{2.4cm}@{}}
\toprule
\textbf{Nutzen} & \textbf{Fälle} & \textbf{Anteil}\\
\midrule
direkt (Geimpfte) & $44\,650$ & $63\,\%$\\
indirekt (Ungeimpfte) & $26\,100$ & $\mathbf{37\,\%}$\\
\midrule
Summe & $70\,750$ & $100\,\%$\\
\bottomrule
\end{tabular}
\end{center}

\schritt{Schritt 5 --- Die Deutung}
\begin{center}
\begin{minipage}{0.9\textwidth}
\textbf{Mehr als ein Drittel des Nutzens fällt bei Menschen an, die sich nicht haben impfen lassen.}
\end{minipage}
\end{center}

\schritt{Schritt 6 --- Je Geimpftem gerechnet}
\[ \frac{70\,750}{50\,000} = 1{,}42 \ \text{vermiedene Fälle je Impfung} \]

Eine Impfung verhindert im Mittel mehr als eine Erkrankung --- sie schützt die geimpfte Person \kursiv{und} einen halben weiteren Menschen.

\begin{center}
\begin{minipage}{0.88\textwidth}
\kursiv{Das ist der rechnerische Kern der Debatte um Impfpflichten. Die Entscheidung, sich nicht impfen zu lassen, betrifft nicht nur die entscheidende Person --- sie entzieht anderen einen Schutz, den diese nicht selbst herstellen können. Das Modell liefert dazu die Zahlen; die Abwägung liefert es nicht.}
\end{minipage}
\end{center}

\schritt{Schritt 7 --- Die Abhängigkeit von der Quote}
\begin{center}
\renewcommand{\arraystretch}{1.45}
\begin{tabular}{@{}P{1.8cm}P{2.6cm}P{2.6cm}@{}}
\toprule
$p$ & \textbf{vermieden je Impfung} & \\
\midrule
$0{,}10$ & $1{,}25$ & \\
$0{,}30$ & $1{,}31$ & \\
$0{,}50$ & $1{,}42$ & \\
$0{,}59$ & $1{,}51$ & \textcolor{gruen}{knapp unter der Schwelle}\\
\bottomrule
\end{tabular}
\end{center}

Je näher an der Schwelle, desto mehr bewirkt jede einzelne Impfung --- die Nichtlinearität aus Abschnitt 8.3 in einer anderen Gestalt.
\end{aufgabe}

\begin{geschichte}
\boxtitelg{Die einzige ausgerottete Krankheit des Menschen}
\Person{Edward Jenner} zeigte 1796, dass eine Infektion mit Kuhpocken gegen Menschenpocken schützt. Das Wort \kursiv{Vakzine} stammt vom lateinischen \kursiv{vacca} --- Kuh.

\begin{center}
\begin{minipage}{0.9\textwidth}
\kursiv{Jenner war nicht der Erste. Die Beobachtung, dass Melkerinnen keine Pocken bekamen, war unter Landärzten bekannt; der Bauer Benjamin Jesty hatte schon 1774 seine Familie geimpft. Jenners Verdienst war, es systematisch zu prüfen und zu veröffentlichen.}
\end{minipage}
\end{center}

\schritt{Die Ausrottung}
Die Weltgesundheitsorganisation beschloss 1967 ein weltweites Ausrottungsprogramm. Die Pocken hatten damals noch etwa zwei Millionen Todesfälle im Jahr zur Folge.

Der letzte natürlich erworbene Fall trat 1977 in Somalia auf; 1980 wurden die Pocken für ausgerottet erklärt.

\begin{center}
\begin{minipage}{0.9\textwidth}
\kursiv{Entscheidend war nicht die flächendeckende Impfung, sondern ein Strategiewechsel: Statt möglichst viele zu impfen, suchte man jeden Fall auf und impfte sein gesamtes Umfeld. Diese \Begriff{Riegelungsimpfung} erreichte die Ausrottung mit Quoten, die deutlich unter der Herdenschwelle lagen.}
\end{minipage}
\end{center}

\schritt{Warum das mathematisch funktioniert}
Die Schwellenformel setzt gleichmäßige Durchmischung voraus. Die Riegelungsimpfung nutzt das Gegenteil aus: Sie impft genau dort, wo das Ansteckungsgeschehen stattfindet.

\begin{center}
\begin{minipage}{0.9\textwidth}
\kursiv{Damit ist sie ein Beispiel dafür, dass eine Bruchstelle des Modells --- die verletzte Durchmischungsannahme --- sich auch \kursiv{nutzen} lässt. Vier Eigenschaften der Pocken machten es möglich: kein Tierreservoir, keine unbemerkten Verläufe, ein wirksamer und hitzestabiler Impfstoff, und ein Erreger ohne relevante Veränderlichkeit.}
\end{minipage}
\end{center}

\schritt{Warum es bei anderen nicht gelingt}
\begin{center}
{\small
\renewcommand{\arraystretch}{1.45}
\begin{tabular}{@{}P{2.2cm}P{2.6cm}P{3.2cm}@{}}
\toprule
\textbf{Erreger} & \textbf{Hindernis} & \\
\midrule
Polio & unbemerkte Verläufe & über $99\,\%$ ohne Lähmung\\
Masern & sehr hohes $\Rnull$ & Schwelle bei $94\,\%$\\
Grippe & Veränderlichkeit & jährlich neuer Impfstoff\\
Gelbfieber & Tierreservoir & Affen als Rückzugsort\\
\bottomrule
\end{tabular}}
\end{center}
\end{geschichte}

\begin{forschung}
\boxtitelp{Wohin sich die Impfmodellierung bewegt}
\begin{center}
{\small
\renewcommand{\arraystretch}{1.45}
\begin{tabular}{@{}P{2.8cm}P{5.4cm}@{}}
\toprule
\textbf{Alters\-strukturierte Modelle} & Kontaktmatrizen aus Befragungen; sie zeigen, dass Kinder untereinander weit mehr Kontakte haben als mit Erwachsenen.\\
\textbf{Räumliche Bündelung} & Modelle mit örtlich schwankender Impfquote; sie sagen Ausbrüche vorher, die Durchschnittsmodelle ausschließen.\\
\textbf{Verhaltens\-modelle} & Die Impfbereitschaft hängt von der wahrgenommenen Gefahr ab --- und die sinkt, wenn die Impfung wirkt.\\
\bottomrule
\end{tabular}}
\end{center}

\schritt{Die dritte Zeile}
Sie beschreibt eine Rückkopplung, die dem Modell dieses Kapitels fehlt: Eine erfolgreiche Impfkampagne lässt die Krankheit verschwinden --- und damit auch die Erinnerung an ihre Schwere.

\begin{center}
\begin{minipage}{0.88\textwidth}
\kursiv{Man kann daraus ein Modell bauen, in dem die Impfquote selbst eine Zustandsgröße ist. Solche Modelle zeigen regelmäßig Schwingungen: Die Quote steigt nach Ausbrüchen, sinkt in ruhigen Zeiten, bis ein neuer Ausbruch folgt. Das Muster entspricht dem, was man bei Masern in mehreren Ländern beobachtet hat --- und es ist wieder die verzögerte Rückkopplung aus Kapitel~3.}
\end{minipage}
\end{center}

\schritt{Die Kontaktmatrizen}
Groß angelegte Befragungen zum Kontaktverhalten haben die Modellierung verändert: Statt einer Durchmischungsannahme benutzt man gemessene Kontaktzahlen zwischen Altersgruppen.

\begin{center}
\begin{minipage}{0.88\textwidth}
\kursiv{Das Ergebnis ist eine Matrix mit starkem Übergewicht auf der Diagonale --- Menschen treffen überwiegend Gleichaltrige --- und einem Nebenband, das Eltern und Kinder verbindet. Aus ihr ergibt sich $\Rnull$ als größter Eigenwert, wieder mit dem Werkzeug aus Kapitel~6.}
\end{minipage}
\end{center}
\end{forschung}

\begin{merksatz}
\boxtitel{Was Sie jetzt können}
\begin{center}
\renewcommand{\arraystretch}{1.5}
\begin{tabular}{@{}P{3.4cm}P{5.2cm}@{}}
\toprule
\textbf{Satz 8.1} & Herdenschwelle $p_c = 1-1/\Rnull$\\
\textbf{Mit Wirksamkeit} & $p > p_c/\varepsilon$; unmöglich, falls $\varepsilon < p_c$\\
\textbf{Unterhalb der Schwelle} & jede Impfung nützt, zunehmend\\
\textbf{Satz 8.2} & $A \approx L/(\Rnull-1)$\\
\textbf{Alters\-verschiebung} & halbherziges Impfen kann schaden\\
\bottomrule
\end{tabular}
\end{center}

\begin{center}
\begin{minipage}{0.88\textwidth}
\kursiv{Im nächsten Kapitel geben wir die Durchmischungsannahme auf --- und stellen Kontakte als Netzwerk dar. Dabei zeigt sich, dass die Schwelle bei manchen Netzen überhaupt nicht existiert.}
\end{minipage}
\end{center}
\end{merksatz}
